\section{Related Work}





% \verb|E-Bench| addresses these limitations by decoupling automatic environment and task construction: each domain is a reusable, controllable, database-backed product world shared across many closed-loop tasks. LLM agents then synthesize tasks by inspecting this environment, executing intended state changes, and recording verified database diffs as trajectory-agnostic ground truth, enabling deterministic grading without task-specific engineering or LLM-based evaluation.




Interaction with speech language models has gradually evolved from conventional turn based human computer dialogue toward full duplex interaction~\cite{surveyduplex} and agentic task execution. Along this trajectory, existing research can be broadly organized into three closely related directions: audio interaction models, Omni interaction models, and voice agent systems for practical task execution. Although these lines of research pursue different objectives, they share a common goal: moving beyond the conventional paradigm in which a model waits for a complete user utterance before responding, toward agents that can continuously perceive their environment, infer the current interaction state, and engage in communication in a natural and timely manner.

\subsection{Speech Language Models and Full Duplex Interaction} Early speech language models~\cite{ji2024wavchat,an2024funaudiollm,kimiaudio,stepaudio1} achieved substantial progress in speech understanding~\cite{qwenaudio,qwen2audio} and speech generation~\cite{cosyvoice3,controlspeech,ji2024mobilespeech,ji2025wavtokenizer,languagecodec}, but their interaction protocols were still largely governed by explicit dialogue turns. In typical systems, external modules such as Voice Activity Detection (VAD)~\cite{SileroVAD,xu2026fireredasr2s} are used to determine when the user begins and ends an utterance, after which the resulting speech segment is passed to the model for understanding and response generation. Consequently, the model itself has limited control over fundamental interaction decisions. Although this paradigm is effective for conventional question and answer interactions, it becomes less suitable for natural conversations involving hesitation, pauses, interruption, overlapping speech, and background interference. Recent work has therefore increasingly focused on enabling speech models to directly model interaction timing and conversational state.

BayLing-Duplex~\cite{BayLing-Duplex} incorporates decisions about when to listen, when to speak, and when to terminate the current response directly into a single autoregressive model. By introducing a small number of dedicated state tokens, the model is able to make interaction decisions during streaming speech generation without relying on an additional turn taking controller. This design treats interaction timing as part of the model's autoregressive prediction process rather than as an external system component. Qwen-Audio-3.0-Realtime~\cite{qwen_audio_realtime_2026} explores realtime speech interaction from a streaming perspective. Instead of treating each utterance as a complete acoustic segment, the continuous audio stream is processed incrementally \textbf{in chunks}, allowing the model to continuously acquire contextual information and generate responses with reduced interaction latency. Audio Interaction Model~\cite{AudioInteractionModel} extends the scope of realtime interaction from conversational speech to general online audio interaction. The Model continuously listens to environmental sounds and user instructions and determines whether a response is necessary. The emphasis on proactive interaction is particularly important because the model is no longer required to respond only after an explicitly defined user turn. Instead, it can initiate a response when the semantics of the ongoing audio stream indicate that intervention is appropriate. Seeduplex~\cite{seeduplex2026} further investigates the challenges of continuous listening in realistic acoustic environments. Its focus goes beyond enabling simultaneous listening and speaking to include attentive listening and robust interference suppression. In particular, the model is required to distinguish relevant user speech from background sounds and speech produced by other speakers while maintaining the ability to respond appropriately to user interruptions. These capabilities highlight an important aspect of full duplex interaction: natural interaction depends not only on simultaneous input and output, \textbf{but also on the model's ability to maintain an appropriate interaction state under complex acoustic conditions.} GPT-Live~\cite{openai_gpt_live_system_card_2026} represents another important line of development toward natural realtime speech interaction. GPT-Live adopts a full-duplex architecture similar to Moshi~\cite{defossez2024moshi}, allowing it to continuously process incoming audio and generate speech output in parallel. This allows the system to make interaction decisions during an ongoing conversation, including whether to continue listening, respond, pause, or interrupt. 

These studies have substantially advanced speech interaction beyond the conventional turn based paradigm. The central question has shifted from whether a model can understand and generate speech to whether it can continuously listen, speak, and regulate its own participation in a conversation. Nevertheless, most of these efforts remain centered on audio based interaction. Their understanding of the surrounding environment is therefore primarily derived from acoustic information, and their interaction capabilities are still largely evaluated within relatively constrained conversational settings.

\subsection{Omni Modalities and Continuous Multimodal Interaction} The integration of visual, audio, and video information into a unified foundation model has opened a new research direction toward Omni interaction. Compared with speech only systems, Omni models have access to both auditory and visual context, enabling the model to reason about not only what the user says, but also the surrounding scene and its temporal evolution. This additional context provides new opportunities for interaction, particularly in situations where the meaning of speech depends on visual information. Representative models such as Qwen Omni~\cite{qwen3.5omni,qwen3omni} series integrate multiple input modalities within a unified foundation model and support streaming speech generation. These systems provide a general foundation for realtime multimodal interaction by enabling the model to jointly interpret linguistic, acoustic, visual, and temporal information within a shared context.



MiniCPM-o 4.5~\cite{cui2026minicpm} further explores continuous multimodal interaction through Omni Flow. Rather than treating multimodal perception and response generation as independent stages, Omni Flow places multimodal inputs and outputs on a unified temporal axis, allowing the model to continuously perceive visual and audio signals while generating responses. This design is particularly relevant to realtime interaction because the model can maintain a continuous temporal representation of the interaction rather than repeatedly resetting its context at individual turns. MiniCPM-o 4.5 demonstrates proactive behaviors based on continuously observed environmental information. JoyAI-VL-Interaction~\cite{yao2026joyai} focuses on interaction over continuous visual streams. Instead of treating vision as a static source of contextual information, it considers continuous visual perception as part of the interaction process itself. The model observes changes in the environment over time and determines whether these changes warrant a response or further interaction. SeedRealtime~\cite{seedrealtime2026} further advances this direction by introducing native audio visual full duplex interaction. It jointly integrates audio, video, and text within a unified architecture and enables continuous interaction over multimodal streams. Beyond simply combining multiple input modalities, SeedRealtime demonstrates that audio visual integration can directly improve interaction quality. For example, visual context can help resolve phonetic ambiguity, interpret temporal references, and connect what is being seen, heard, and said within the same interaction process. The model also demonstrates proactive interaction by responding to changes in the observed environment without requiring an explicit user request.


% These developments indicate that Omni interaction introduces capabilities that are difficult to obtain from audio alone. By jointly modeling visual, acoustic, and temporal information, a model can use scene context to resolve ambiguous speech, interpret references, track changes in the environment, and infer user intent more accurately. More importantly, the model can make interaction decisions based not only on the user's explicit speech, but also on the state of the surrounding environment.

\subsection{From Omni Conversation to Voice Agents}

Although the above approaches have improved the naturalness of conversational interaction, practical deployment introduces a further challenge. Users increasingly expect AI systems not only to converse with them, but also to perform meaningful tasks, such as writing code~\cite{glm5}, accessing files, invoking external tools~\cite{voxmind}, and executing multi step workflows. In these settings, conversational interaction and task execution are closely coupled, yet they are still commonly implemented as partially independent capabilities.

Qwen Audio Agent~\cite{qwen_audio_agent_2026} represents a system oriented toward this problem. It provides a realtime voice runtime that enables an agent to maintain continuous spoken interaction while delegating more complex tasks to a backend workflow. Such tasks can include code generation, computer file operations, and other work related activities. This design moves voice interaction beyond a conversational interface and toward a persistent interface to an agent workflow. GPT-Live with CodeX~\footnote{\url{https://help.openai.com/en/articles/20001274-chatgpt-voice?utm_source=chatgpt.com}} and Claude Voice Mode~\footnote{\url{https://support.claude.com/en/articles/11101966-use-voice-mode?utm_source=chatgpt.com}} demonstrate product oriented approaches in which natural voice interaction is integrated with general purpose AI assistant capabilities. These systems make increasingly complex AI functionality accessible through continuous spoken interaction, thereby reducing the distinction between a conversational interface and a general AI assistant.

Despite these advances, current voice agent systems still face substantial challenges in realistic interactive workflows. Throughout task execution, users may interrupt the agent and provide new constraints, ask follow up questions, or correct previously stated information. The agent may also need to proactively communicate intermediate results, solicit clarification, and request additional details to ensure that the task proceeds correctly. A brief backchannel does not necessarily indicate that the current task should terminate. In multi party environments, the system may need to distinguish relevant speech from other speakers while simultaneously handling background noise and changes in the visual environment. These scenarios require the model to reason jointly about interaction state, environmental state, and task state. Treating speech only as an input and output interface is therefore insufficient for truly natural agent interaction.

Motivated by these developments, we propose \textbf{an omni interaction agent called Gander} which studies a more general form of interaction in which a single end to end omni model simultaneously supports natural realtime conversation and agentic task execution. The key question is whether conversational interaction, multimodal perception, and task execution can be integrated into a single continuous process rather than being handled as independent system components. Gander aims to bring these capabilities together within a unified Omni Interaction Agent. Rather than focusing solely on conversational interaction, Gander is designed to move naturally between open ended conversation and practical task execution. Rather than treating speech as merely an interface to an external agent, Gander incorporates interaction state into the end to end modeling process, allowing conversational behavior and task execution to be coordinated within the same ongoing interaction.
